\begin{textsample}{POS Dim 2 – to – Score 22.00 – to/to\_em\_16.txt}  \label{ex:f2_pos_069}
2.
Objetivos Os \textbf{Itinerários} Formativos têm como função precípua a ampliação das aprendizagens nas diversas áreas do conhecimento, de modo a garantir o desenvolvimento das competências gerais da BNCC, as habilidades dos \textbf{itinerários} formativos associadas às competências gerais e as habilidades específicas dos \textbf{itinerários} formativos associadas aos eixos estruturantes.
Desse modo, os objetivos dos \textbf{Itinerários} Formativos, conforme estabelece a \textbf{Portaria} MEC n{\EmojiFont °} 1.432/2018, são : a ) Aprofundar as aprendizagens relacionadas às competências gerais, às Áreas de Conhecimento e/ou à Formação \textbf{Técnica} e Profissional ; b ) Consolidar a formação integral dos estudantes, desenvolvendo a autonomia necessária para que realizem seus projetos de vida ; c ) Promover a incorporação de valores universais, como ética, liberdade, democracia, justiça social, pluralidade, solidariedade e sustentabilidade ; d ) Desenvolver habilidades que permitam aos estudantes ter uma visão de mundo ampla e heterogênea, tomar decisões e agir nas mais diversas situações, seja na escola, seja no trabalho, seja na vida.
Nesse sentido, o \textbf{artigo} 12 da Resolução CNE n{\EmojiFont °} 03, de 21 de \textbf{novembro} de 2018, que \textbf{atualiza} as \textbf{Diretrizes} Curriculares Nacionais para o Ensino Médio, estabelece que a organização dos \textbf{Itinerários} deve considerar : Art. 12.
A partir das áreas do conhecimento e da formação \textbf{técnica} e profissional, os \textbf{itinerários} formativos devem ser organizados, considerando : I - linguagens e suas tecnologias : aprofundamento de conhecimentos estruturantes para aplicação de diferentes linguagens em contextos sociais e de trabalho, estruturando arranjos curriculares que permitam estudos em línguas vernáculas, estrangeiras, clássicas e indígenas, Língua Brasileira de Sinais ( LIBRAS ), das artes, design, linguagens digitais, corporeidade, artes cênicas, roteiros, produções literárias, dentre outros, considerando o contexto local e as possibilidades de \textbf{oferta} pelos sistemas de ensino ; II - matemática e suas tecnologias : aprofundamento de conhecimentos estruturantes para aplicação de diferentes conceitos matemáticos em contextos sociais e de trabalho, estruturando arranjos curriculares que permitam estudos em resolução de problemas e análises complexas, funcionais e não-lineares, análise de dados estatísticos e probabilidade, geometria e topologia, robótica, automação, inteligência artificial, programação, jogos digitais, sistemas dinâmicos, dentre outros, considerando o contexto local e as possibilidades de \textbf{oferta} pelos sistemas de ensino ; III - ciências da natureza e suas tecnologias : aprofundamento de conhecimentos estruturantes para aplicação de diferentes conceitos em contextos sociais e de trabalho, organizando arranjos curriculares que permitam estudos em astronomia, metodologia, física geral, clássica, molecular, quântica e mecânica, instrumentação, ótica, acústica, química dos produtos naturais, análise de fenômenos físicos e químicos, meteorologia e climatologia, microbiologia, imunologia e parasitologia, ecologia, nutrição, zoologia, dentre outros, considerando o contexto local e as possibilidades de \textbf{oferta} pelos sistemas de ensino ; IV - ciências humanas e sociais aplicadas : aprofundamento de conhecimentos estruturantes para aplicação de diferentes conceitos em contextos sociais e de trabalho, estruturando arranjos curriculares que permitam estudos em relações sociais, modelos econômicos, processos políticos, pluralidade cultural, historicidade do universo, do homem e natureza, dentre outros, considerando o contexto local e as possibilidades de \textbf{oferta} pelos sistemas de ensino ; V - formação \textbf{técnica} e profissional : desenvolvimento de programas educacionais inovadores e \textbf{atualizados} que promovam efetivamente a \textbf{qualificação} profissional dos estudantes para o mundo do trabalho, objetivando sua \textbf{habilitação} profissional tanto para o desenvolvimento de vida e \textbf{carreira}, quanto para adaptar -se às novas condições ocupacionais e às exigências do mundo do trabalho contemporâneo e suas contínuas transformações, em condições de competitividade, produtividade e inovação, considerando o contexto local e as possibilidades de \textbf{oferta} pelos sistemas de ensino.
Além disso, para a definição e elaboração dos \textbf{Itinerários} Formativos, deve se considerar que o \textbf{currículo} escolar mais \textbf{flexível} busca responder à diversidade de aptidões, referência e objetivos dos estudantes, bem como, dialogar com os diversos contextos educacionais, sociais, culturais e produtivos do país.
Em razão disso, alguns pontos devem ser observados : a ) oportunizar aos estudantes, mesmo em contexto com reduzida \textbf{oferta} de \textbf{itinerários}, a possibilidade de fazer escolhas que coadunam com seu interesse.
Nesse caso, a rede de ensino deve \textbf{viabilizar} condições para \textbf{atender} aos anseios dos estudantes, por meio de arranjos curriculares na forma de unidades \textbf{eletivas}, projetos, atividades integradoras, com organização multisseriada, ou ainda, por meio da \textbf{oferta} não presencial, com uso de recursos tecnológicos ou não. b ) dar autonomia às unidades escolares para propor, com participação dos estudantes, a escolha dos \textbf{itinerários}, e posteriormente, permitir que elas, tendo por base a participação democrática, o contexto socio-histórico-cultural, o meio social, definam e construam suas unidades curriculares e seus \textbf{Itinerários} Formativos, a partir das \textbf{diretrizes} gerais expressas na BNCC, do Sistema de Ensino e do \textbf{Currículo} Estadual. c ) possibilitar aos estudantes, o aproveitamento das aprendizagens realizadas em \textbf{instituições} parceiras, mediante regulação do Sistema de Ensino, de modo a fortalecer a autonomia e o protagonismo dos estudantes.
Nesse sentido, as unidades escolares devem aproveitar os estudos e os conhecimentos no ensino formal e a experiência extraescolar, conforme preconiza as DCNEM. d ) definir o perfil de saída dos estudantes por ocasião da conclusão dos \textbf{Itinerários} Formativos, com o objetivo de conduzir os estudantes ao aprofundamento em uma ou mais áreas do conhecimento ou em campo de formação profissional, de modo a delinear o \textbf{prosseguimento} dos estudos e a preparação para o mundo do trabalho. e ) coadunar os \textbf{Itinerários} Formativos com o \textbf{disposto} na \textbf{Portaria} MEC n{\EmojiFont °} 1.432, de 28 de dezembro de 2018, que estabelece os referenciais para elaboração dos \textbf{itinerários} formativos conforme \textbf{preveem} as \textbf{Diretrizes} Nacionais do Ensino Médio.
Sendo assim, devem ser organizados a partir dos eixos estruturantes : Investigação Científica ; Processos Criativos ; Mediação e Intervenção Sociocultural ; e Empreendedorismo, conforme as tabelas 1 e 2.

% matched lemmas: artigo, atender, atualizar, carreira, currículo, diretriz, disposto, eletivo, flexível, habilitação, instituição, itinerário, novembro, oferta, portaria, prever, prosseguimento, qualificação, técnico, viabilizar
\end{textsample}
